\chapter*{Abstract}
\addcontentsline{toc}{chapter}{Abstract}

LLaVA is a pioneering vision-language model that grounds a large language model (LLM) in images to enable multimodal assistance. This thesis investigates how visual instruction tuning reshapes the backbone's \emph{representations geometry}, where along the network's depth these changes concentrate, and how they relate to downstream performance. We analyze all backbone layers using layer- and head-wise (within attention layers) representation similarity, profile the intrinsic dimension, and evaluate the impact of layer patching. Our results show that the effects of visual instruction tuning localize to the early-mid layers, where the similarity between the pretrained and fine-tuned models drops the most, especially for multimodal inputs. The same region shows strong geometric compression (lower intrinsic dimension), and cumulatively patching early-mid layers from the fine-tuned checkpoint into the pretrained one recovers most gains on visual question answering and image captioning, whereas patching late layers yields diminishing returns. Overall, visual instruction tuning teaches LLaVA to \emph{see} by primarily modifying the early-mid region of the network, while late layers act as largely modality-agnostic decoders---suggesting targeted, cost-effective fine-tuning strategies and clarifying how multimodal LLMs integrate modalities. The code for this thesis is available at \url{https://github.com/luispky/LLaVA_Representations_Analysis.git}.

\chapter*{Abstract (Italian)}
\addcontentsline{toc}{chapter}{Abstract (Italian)}

LLaVA è un modello di visione e linguaggio pionieristico, che collega un \textit{large language model (LLM)} alle immagini per abilitare l'assistenza multimodale. Questa tesi analizza come il \textit{visual instruction tuning} rimodella la geometria delle rappresentazioni del modello, in che punti della rete si concentrano questi cambiamenti e in che modo questi si riflettono sulle prestazioni a valle. Analizziamo tutti i \textit{layer} del modello utilizzando misure di similarità tra rappresentazioni a livello di layer e di \textit{head} (all'interno dei layer di attenzione), oltre a misurarne la dimensione intrinseca e valutare la possibilità di trasferire interi layer. I nostri risultati mostrano che gli effetti del visual instruction tuning si localizzano principalmente nella prima metà del modello, dove la similarità tra il modello originale e quello \textit{fine-tuned} è più bassa, soprattutto per input multi-modali. La stessa regione mostra una forte compressione geometrica (dimensione intrinseca più bassa) e il trasferimento cumulativo dei primi layer dal modello fine-tuned in quello originale consente di recuperare gran parte dei miglioramenti in \textit{visual question answering} e \textit{image captioning}, mentre il trasferimento degli strati finali comporta rendimenti decrescenti. In sintesi, il visual instruction tuning insegna a LLaVA a vedere modificando principalmente la porzione iniziale della rete, mentre gli strati successivi agiscono in gran parte come decoder indipendenti dalla modalità---suggerendo strategie di fine-tuning mirate ed economicamente efficienti e chiarendo come gli LLM multi-modali integrino le diverse modalità. Il codice sviluppato per questa tesi è disponibile all'indirizzo \url{https://github.com/luispky/LLaVA_Representations_Analysis.git}.